\begin{textsample}{POS Dim 1 – ms – Score 35.00 – ms/ms\_em\_29.txt}  \label{ex:f1_pos_203}
3.1 Área de Conhecimento : Linguagens e suas Tecnologias A Base Nacional Comum Curricular ( BNCC ) determina \textbf{que} a área de Linguagens e suas Tecnologias explore os diversos tipos de linguagem, expandindo as capacidades de expressões artísticas, corporais e linguísticas dos estudantes.
A participação dos \textbf{discentes} nas diferentes práticas de linguagem colabora para sua formação integral, \textbf{que} não se deve limitar ao aperfeiçoamento de competências cognitivas, mas também das socioemocionais, a fim de \textbf{que} a aprendizagem auxilie na resolução de problemas e produção de conhecimentos \textbf{que} reflitam no desenvolvimento da sociedade.
A educação integral é \textbf{um} dos princípios específicos da etapa do Ensino Médio e todas as suas modalidades, conforme determina o artigo 5º, inciso I da Resolução n. 3, de 21 de novembro de 2018, \textbf{que} atualiza as Diretrizes Nacionais para o Ensino Médio ( DCNEM ).
Assim, o \textbf{que} se pretende é \textbf{que} a escola prepare o estudante para aprender continuamente, aplicar os conhecimentos construídos em seu benefício e de \textbf{uma} coletividade.
No artigo 6, inciso I, define -se : > [... ] formação integral : é o desenvolvimento intencional dos aspectos físicos, cognitivos e socioemocionais do estudante por meio de processos educativos significativos \textbf{que} promovam a autonomia, o comportamento cidadão e o protagonismo na construção de seu projeto de vida ( BRASIL, 2018a ).
Segundo a Base Nacional Comum Curricular, ao desenvolver as competências e habilidades da área de Linguagens e suas Tecnologias, o estudante tem a oportunidade de consolidar e ampliar sua capacidade “ de uso e reflexão sobre as linguagens – artísticas, corporais e verbais ” ( BRASIL, 2018c, p. 481 ), ou seja, não será somente leitor, mas \textbf{autor} de diversos gêneros discursivos \textbf{contemporâneos}, muitas vezes, relacionados ao universo digital.
Cabe ressaltar \textbf{que} a área de Linguagens e suas Tecnologias ( LGG ) propõe \textbf{que} se ultrapasse a linguagem verbal, considerando todos os elementos \textbf{que} contribuem para a significação do texto.
A proposta é de \textbf{uma} escola \textbf{que} \textbf{aproxime} os estudantes das mais diversas produções, em especial, as artísticas e culturais, dentro dos contextos regionais.
Conforme afirma Rojo : > Assim como foi capaz de popularizar os impressos, \textbf{urge} \textbf{que} a escola se \textbf{preocupe} com o acesso a outros espaços valorizados de cultura ( museus, bibliotecas, teatros, espetáculos ) e a outras mídias ( analógicas e digitais ) ( ROJO, 2009, p. 52 ).
Na etapa do Ensino Médio, a BNCC define como prioridade cinco campos de atuação social, \textbf{que} orientam para práticas de linguagem \textbf{que} possibilitem aos estudantes “ vivenciarem experiências significativas em diferentes mídias ( impressa, digital, analógica ) ”.
Os campos de atuação objetivam o avanço da educação escolar ao propor \textbf{que} a cultura digital, os multiletramentos e os novos letramentos sejam acrescidos à cultura do impresso ( escrita ), propiciem a análise de elementos visuais ( imagens estáticas ou em movimento ), sonoros ( músicas, ruídos, sonoridades ), verbais ( oral ou \textbf{visual-motora}, como Libras e escrita ) e corporais ( gestuais, cênicas, dança ), ou seja, considerem diferentes semioses.
Esses campos de atuação compreendem : - campo da vida pessoal : processos de construção de identidade e de Projetos de Vida ; - campo das práticas de estudo e pesquisa : construção do conhecimento científico para aprender a aprender ; - campo jornalístico-midiático : construção de consciência crítica e seletiva em relação à produção e circulação de informações ; - campo de atuação na vida pública : reflexão e participação na vida pública, pautado pela ética ; - campo artístico : reconhecimento, valorização, fruição e produção de manifestações artísticas em geral.
Nessa perspectiva, o Currículo de Referência de Mato Grosso do Sul — Ensino Médio da área de Linguagens e suas Tecnologias — estrutura -se em \textbf{um} Organizador Curricular, \textbf{que} corresponde a \textbf{um} quadro composto por cinco ( 05 ) colunas, contendo : eixo temático ( \textbf{que} em LGG equivale ao campo de atuação social ), habilidades, componente curricular, objeto de conhecimento e sugestões didáticas, os quais apresentam os saberes \textbf{que} devem ser aprofundados com os estudantes, por meio de competências e habilidades, em cada componente curricular da Área de Conhecimento, de determinado ano, dando respaldo pedagógico aos professores.
Considerando \textbf{que} os estudantes possuem ritmos de aprendizagem muito diferentes \textbf{uns} dos outros, as graduações das complexidades das habilidades devem acompanhar o desenvolvimento de cada indivíduo.
O Estado de Mato Grosso do Sul optou por manter os componentes curriculares, já existentes, na formação geral, partindo de sugestões didáticas \textbf{que} desenvolvam as competências e habilidades tanto por Área de Conhecimento quanto interdisciplinarmente.
No espaço do organizador curricular da Formação Geral Básica ( FGB ), destinado às sugestões didáticas, são apresentadas práticas \textbf{metodológicas} \textbf{que} proporcionam aos professores orientações de estudo por área. \textbf{Enquanto} sugestões, é importante destacar \textbf{que} essas ações não devem ser compreendidas como obrigatórias, mas como possibilidades \textbf{que} precisam fazer sentido para as unidades escolares, \textbf{que} podem adaptá -las ou usá -las como inspiração, ou definir suas próprias ações didáticas, considerando as especificidades de cada localidade.
Essas sugestões estão atreladas aos Temas \textbf{Contemporâneos}, à prática da presença pedagógica, ao uso de metodologias ativas, à pesquisa científica, ao protagonismo juvenil e à autoria, assegurando, assim, o desenvolvimento global do estudante e a construção de seu Projeto de Vida, conforme previsto no artigo 11 das DCNEM : > Art. 11 A formação geral básica é composta por competências e habilidades previstas na Base Nacional Comum Curricular ( BNCC ) e articuladas como \textbf{um} todo indissociável, enriquecidas pelo contexto histórico, econômico, social, ambiental, cultural local, do mundo do trabalho e da prática social, e deverá ser organizada por áreas de conhecimento [... ] ( BRASIL, 2018a ).
Assim, as atividades sugeridas em LGG mobilizam conjuntos de competências e habilidades \textbf{que} se \textbf{inter-relacionam}, denominados agrupamentos, possibilitando o aprofundamento do conhecimento, a partir de estratégias \textbf{metodológicas} \textbf{que} se complementam com os componentes curriculares da área com as demais, quando possível.
Destaca -se \textbf{que} cada componente ainda explora suas especificidades, proporcionando ao estudante a capacidade de identificar, analisar e posicionar -se, criticamente, diante de situações da realidade.
No Estado de Mato Grosso do Sul, a área de Linguagens e suas Tecnologias é composta pelos seguintes componentes curriculares : Arte, Educação Física, Língua Inglesa e Língua Portuguesa, na FGB, e Língua Espanhola no Núcleo Integrador¹⁰, para casos em \textbf{que} a escola ofereça \textbf{uma} segunda língua estrangeira, conforme o \textbf{que} define a Lei 13.415, de 2017, em seu artigo 35 A, parágrafo 4º : > Os currículos do ensino médio incluirão, obrigatoriamente, o estudo da língua inglesa e poderão ofertar outras línguas estrangeiras, em caráter optativo, preferencialmente o espanhol, de acordo com a disponibilidade de oferta, locais e horários definidos pelos sistemas de ensino ( BRASIL, 2017 ).
Ressalta -se \textbf{que} essa organização tem como objetivos potencializar a construção do conhecimento e articular os componentes curriculares, partindo de \textbf{um} ponto em comum, nesse caso, das competências e habilidades a serem desenvolvidas, bem como do campo de atuação social.
Para a área de Linguagens e suas Tecnologias apresentam -se dois organizadores curriculares : o primeiro explora o desenvolvimento de habilidades e competências, por área, e o segundo enfatiza as habilidades específicas de Língua Portuguesa, voltadas às práticas de leitura, escrita e oralidade, de diversos gêneros discursivos e ao estudo do texto \textbf{literário}, tendo em vista os mecanismos linguísticos e multissemióticos, visando à formação do “ lautor ” ( ROJO, 2013, p. 20 ). ¹⁰ O Núcleo Integrador compreende os componentes curriculares Projeto de Vida, Empreendedorismo Social, Ciências Integradas e Novas Tecnologias, Matemática Criativa e Resolução de Problemas, Linguagens e Interartes, Língua Espanhola e Intervenção Comunitária \textbf{que} se articulam de forma integrada e \textbf{permeiam} por todas as áreas de conhecimento e possibilitam o desenvolvimento de competências e habilidades específicas da BNCC e dos Itinerários Formativos. ( MATO GROSSO DO SUL, 2020 ).
O desenvolvimento das competências e habilidades foi definido agrupando aquelas \textbf{que} tinham \textbf{uma} relação próxima, de forma \textbf{que} pudessem ser trabalhadas em \textbf{uma} sequência didática \textbf{que} passasse pelos principais níveis de aprendizagem, conforme a Taxonomia de Bloom revisada em 2001, chegando ao sexto nível : o da criação e produção ( figura 1 ).
Figura 1 - Taxonomia de Bloom revisada O processo de seriação foi realizado da seguinte maneira : agrupamento de habilidades \textbf{que} se completam ; definição de objetos de conhecimento ; produção de sugestões didáticas ; análise do grau de dificuldade dos agrupamentos de habilidades por componente ; e alinhamento por área do conhecimento e seriação.
Nesse sentido, todas as sugestões didáticas oferecem ao menos \textbf{uma} forma de materialização do conhecimento, a partir de \textbf{uma} produção \textbf{que} utilize as linguagens verbal, corporal, visual e/ou sonora explorando, assim, diferentes semioses e empregando, em diversas ocasiões, as ferramentas próprias da cultura digital.
No organizador, os componentes curriculares aparecem por ordem alfabética : Arte, Educação Física, Língua Espanhola, Língua Inglesa e Língua Portuguesa.
No componente curricular Arte, etapa do Ensino Médio, conforme a BNCC, o arte educador promove : > [... ] o entrelaçamento de culturas e saberes, possibilitando aos estudantes o acesso e a interação com as distintas manifestações culturais \textbf{populares} presentes na sua comunidade.
O mesmo deve ocorrer com outras manifestações presentes nos centros culturais, museus e outros espaços, de modo a propiciar o exercício da crítica, da apreciação e da fruição de \textbf{exposições}, concertos, apresentações musicais e de dança, filmes, peças de teatro, poemas e obras \textbf{literárias}, entre outros, garantindo o respeito e a valorização das diversas culturas presentes na formação da sociedade brasileira, especialmente as de matrizes indígena e africana ( Brasil, 2018c, p. 483 ).
Nesse contexto, o Currículo considera as peculiaridades da comunidade escolar, prioriza a valorização de produções locais, por meio da arte regional, nacional e mundial, para \textbf{que} ocorra a identificação cultural e a visão crítica de sua existência, respeitando a diversidade das culturas locais do Estado de Mato Grosso do Sul.
O objetivo da Arte na Educação Básica não é formar artistas, mas garantir \textbf{que} conhecimentos produzidos contribuam para a \textbf{contextualização} de saberes e práticas artísticas do estudante, valorizando a arte local, a produção cultural e a construção de seu Projeto de Vida.
Ressalta -se a necessidade de \textbf{um} olhar flexível, autônomo e coerente, \textbf{focado} no protagonismo e na articulação das práticas de : criar, ler, produzir, construir, exteriorizar e refletir sobre formas artísticas, comprometido com “ o aprofundamento na pesquisa e no desenvolvimento de processos de criações autorais nas linguagens das artes visuais, do audiovisual, da dança, do teatro, das artes circenses e da música ” ( BRASIL, 2018c, p. 482 ), conforme Barbosa : > [... ] Através das artes é possível desenvolver a \textbf{percepção} e a imaginação, apreender a realidade do meio ambiente, desenvolver a capacidade crítica, permitindo analisar a realidade percebida e desenvolver a criatividade de maneira a mudar a realidade \textbf{que} foi analisada ( BARBOSA, 1998. p. 16 ).
Para o componente curricular Educação Física, na etapa do Ensino Médio, pretende -se \textbf{que} o estudante aprofunde as experiências do Ensino Fundamental sobre as capacidades e limites do corpo ; compreendendo a importância do estilo de vida saudável e do autocuidado na manutenção da saúde.
Dessa forma, o processo vivenciado será consolidado na última fase da Educação Básica : > [... ] ao final do Ensino Médio, o jovem deverá apresentar \textbf{uma} compreensão aprofundada e sistemática acerca da presença das práticas corporais em sua vida e na sociedade, incluindo os fatores sociais, culturais, ideológicos, econômicos e políticos envolvidos nas práticas e nos discursos \textbf{que} circulam sobre elas ( BRASIL, 2018c, p. 495 ).
Ao pensar nas competências e habilidades relacionadas à aprendizagem da área de Linguagem e suas Tecnologias, no \textbf{que} tange às particularidades do Componente Curricular Educação Física, \textbf{logo} tem -se em \textbf{mente} os movimentos corporais diversos : esportes, caminhadas, corridas, treinamentos e outros, no entanto é possível associá -los ao uso das Tecnologias Digitais da Informação e Comunicação ( TDIC ) ao propor novas estratégias \textbf{metodológicas}.
Os videogames, geralmente tidos como inimigo da vida saudável, podem auxiliar na manutenção da qualidade de vida a partir de jogos \textbf{que} favoreçam a atividade física ( exergames ).
Segundo Silva : > Atualmente os Exergames estão sendo comparados a boas atividades de movimentos e práticas, atividade de \textbf{estímulo} ao prazer de se praticar \textbf{uma} atividade e ao mesmo tempo de possibilidade de aprimoramento do aprendizado de crianças e adolescentes em escolas ( SILVA, 2017, p. 6 ).
Assim, propõem -se atividades \textbf{que} auxiliam o professor a compor \textbf{uma} sequência didática \textbf{que} colabore com a qualidade de vida, sem ignorar os avanços tecnológicos ao desenvolver as competências e habilidades da área.
No \textbf{que} \textbf{diz} respeito à Língua Inglesa, a Lei de Diretrizes e Bases, redação dada pela Lei n. 13.415/2017, tornou obrigatório o seu estudo, na etapa do Ensino Médio, devido ao seu “ caráter global – pela multiplicidade e variedade de usos, usuários e funções na \textbf{contemporaneidade} ” ( BRASIL, 2018c. p. 484 ).
Ao propor a aprendizagem de Língua Inglesa, por área de conhecimento, objetiva -se \textbf{que} os estudantes explorem as diversas situações de uso da língua, especialmente na cultura digital, ampliando sua capacidade discursiva e de reflexão, a partir da \textbf{contextualização} das práticas de linguagem nos diversos campos de atuação : > Trata -se, portanto, de expandir os repertórios linguísticos, multissemióticos e culturais dos estudantes, possibilitando o desenvolvimento de maior consciência e reflexão críticas das funções e usos do inglês na sociedade \textbf{contemporânea} – permitindo, por exemplo, problematizar com maior criticidade os motivos pelos quais ela se tornou \textbf{uma} língua de uso global ” ( BRASIL, 2018c.p. 485 ).
O componente curricular Língua Portuguesa, em consonância com a Lei 13.415, é obrigatório nos três anos da etapa do Ensino Médio, juntamente com a Matemática.
Pretende -se, nessa fase, aprofundar os conhecimentos adquiridos na Etapa do Ensino Fundamental referentes a linguagem ( uso, finalidade, significação ) e à leitura do texto \textbf{literário}.
Nesse sentido, a Língua Portuguesa, indissociável da Literatura, busca relacionar -se com as mais diversas linguagens, presentes em diferentes gêneros discursivos, a fim de \textbf{que} o estudante tenha condições de participar, significativamente, das práticas sociais \textbf{contemporâneas}.
Pensando nos níveis de alfabetismo definidos pelo Indicador Nacional de Alfabetismo \textbf{Funcional} — INAF, fica claro \textbf{que} é necessário, ao concluir a Etapa do Ensino Médio, \textbf{que} o estudante tenha “ capacidade de ler textos longos, orientando -se por subtítulo, localizando mais de \textbf{uma} informação, de acordo com as condições estabelecidas, relacionando partes de \textbf{um} texto, comparando dois textos, realizando inferências e sínteses ” ( ROJO, 2009, p.47 ).
Ao construir as sugestões didáticas em Língua Portuguesa, considerou -se o desenvolvimento das práticas de linguagens \textbf{que}, na etapa do Ensino Médio, são chamadas de eixos de integração — leitura, produção de texto, oralidade ( escuta e produção oral ) e análise linguística/semiótica.
Faz -se necessário \textbf{lembrar} \textbf{que}, em alguns momentos, \textbf{um} ou mais eixos de integração estará/estarão em maior evidência, devido ao objeto de conhecimento \textbf{que} será tratado, ou seja, \textbf{ora} enfatiza -se a leitura, \textbf{ora} a oralidade, \textbf{ora} a produção textual.
Com o objetivo de continuar a formação do leitor e ampliar o contato dos estudantes com a Literatura, a Língua Portuguesa explora a riqueza proporcionada por análises contextualizadas, principalmente das obras clássicas, com momentos de investigação, \textbf{correlação} de saberes, resultando no pensamento crítico sobre a realidade.
O estudo da Literatura deve \textbf{supor} o entendimento da linguagem \textbf{literária} \textbf{enquanto} construção linguística \textbf{que} tem sua especificidade histórica, cultural e artística, e deve ser entendido e articulado com as formas de conhecimento e com as condições históricas de sua produção.
Assim, o texto \textbf{literário} é compreendido como \textbf{instância} \textbf{que} reúne \textbf{uma} realidade estética, \textbf{uma} forma de conhecimento \textbf{que} tem como finalidade entender o ser, a existência, o mundo e sua produção de vida nas variadas dimensões de figuração da vida social.
É importante compreender \textbf{que} a literatura faz parte da formação humana e sempre esteve presente no imaginário de cada pessoa.
Segundo Cândido, quando se trata de literatura : > Não há povo e não há \textbf{homem} \textbf{que} possa \textbf{viver} sem ela, isto é, sem a possibilidade de entrar em contacto com alguma espécie de fabulação.
Assim como todos sonham todas as noites, ninguém é capaz de passar as vinte e quatro horas do dia sem alguns momentos de entrega ao universo fabulado ( CÂNDIDO, 1995. p.174 ).
Já em relação ao componente Língua Espanhola, esse encontra -se próximo à realidade cultural e econômica do Estado de Mato Grosso do Sul, devido a sua localização geográfica e suas \textbf{fronteiras} com a Bolívia e o Paraguai.
Para as escolas \textbf{que} optarem pela oferta de \textbf{uma} segunda língua estrangeira, este Currículo de Referência de Mato Grosso do Sul propõe estratégias \textbf{metodológicas} alinhadas aos demais componentes curriculares da área de Linguagens e suas Tecnologias, o \textbf{que} viabiliza \textbf{uma} ação integrada e valoriza essa língua \textbf{que} \textbf{influencia}, fortemente, a cultura regional.
A área de Linguagens e suas Tecnologias contempla 7 ( sete ) competências específicas e 28 ( vinte e oito ) habilidades, apresentando, ainda, 54 ( cinquenta e quatro ) habilidades de Língua Portuguesa, consideradas \textbf{desdobramentos}, o \textbf{que} possibilita o planejamento por área do conhecimento.
Ressalta -se \textbf{que} foram criadas mais duas habilidades, \textbf{uma} para Linguagens e suas Tecnologias ( MS.EM13LGG6.n.01 ) e outra para Língua Portuguesa ( MS.EM13LP4.n.01 ).
As competências e habilidades estão organizadas considerando a proximidade existente entre elas, de formas \textbf{que} possam ser abordadas em \textbf{uma} situação didática \textbf{que} contemple os principais níveis de aprendizagem, conforme a Taxonomia de Bloom, chegando ao sexto nível : o da criação, produção.
O processo de seriação foi realizado da seguinte maneira : 1. agrupamento de habilidades \textbf{que} se completam ; 2. definição de objetos de conhecimento ; 3. produção de sugestões didáticas ; 4. análise do grau de dificuldade dos agrupamentos de habilidades por componente ; 5. alinhamento por área do conhecimento e seriação.
Dessa forma, todas as sugestões didáticas oferecem ao menos \textbf{uma} forma de materialização do conhecimento, a partir de \textbf{uma} produção \textbf{que} explore as linguagens verbal, corporal, visual e/ou sonora, e diferentes semioses, com o emprego das ferramentas próprias da cultura digital.

% matched lemmas: aproximar, autor, contemporaneidade, contemporâneo, contextualização, correlação, desdobramento, discente, dizer, enquanto, estímulo, exposição, focar, fronteira, funcional, homem, influenciar, instância, inter-relacionar, lembrar, literário, logo, mente, metodológico, ora, percepção, permear, popular, preocupar, que, supor, um, urgir, visual-motor, viver
\end{textsample}
